\chapter{算法基础习题}

\section{课堂练习}\qanswerloc{10}

\begin{qitems}[showanalysis]

    \begin{bbox}
        \qitem 算法必须具备输入、输出和 \blankbox 。
        \fourchoices{确定性}{有穷性}{可行性}{以上都是}
        \begin{analysis}
            算法必须具备有穷性，即在有限步后终止。确定性和可行性也是好算法的特征，但有穷性是必须的。
        \end{analysis}
    \end{bbox}

    \begin{bbox}
        \qitem 下列函数的时间复杂度最小的是 \blankbox 。\textwater
        \fourchoices
            {$T_1(n)=n\log_2 n + 5000n$}
            {$T_2(n)=n^2 - 800n$}
            {$T_3(n)=n\log_2 n - 6000n$}
            {$T_4(n)=20000\log_2 n$}
        \begin{analysis}
            $T_4(n) = 20000\log_2 n$ 为对数阶，时间复杂度最小。
        \end{analysis}
    \end{bbox}

    \begin{bbox}
        \qitem 若某算法空间复杂度为 $O(1)$，则表示该算法 \blankbox 。
        \fourchoices
            {不需要任何辅助空间}
            {所需辅助空间大小与问题规模 $n$ 无关}
            {不需要任何空间}
            {所需空间大小与问题规模 $n$ 无关}
        \begin{analysis}
            $O(1)$ 表示所需辅助空间大小与问题规模无关（固定值），但仍可能需要少量辅助空间。
        \end{analysis}
    \end{bbox}

\end{qitems}

\section{课后作业}\qanswerloc{20}

\begin{qitems}[prefix=（, suffix=）, optprefix=, optsuffix=]

    \begin{bbox}
        \qitem 分析以下程序段的时间复杂度。
        \begin{lstlisting}
int func(int n) {
    int i = 0, sum = 0;
    while (sum < n) sum += ++i;
    return i;
}
        \end{lstlisting}
        \fourchoices
            {$O(\log n)$}
            {$O(n^{1/2})$}
            {$O(n)$}
            {$O(n\log n)$}

        \begin{analysis}
            设循环次数为 $k$，则 $\sum_{i=1}^{k} i = \frac{k(k+1)}{2} < n$，
            所以 $k \approx \sqrt{2n} = O(\sqrt{n})$。
        \end{analysis}
    \end{bbox}

    \begin{bbox}
        \qitem 以下关于算法的叙述中，正确的是 \blankbox 。
        \threechoices
            {算法必须有输入和输出}
            {算法必须用流程图表示}
            {算法的时间复杂度越低越好}
    \end{bbox}

\end{qitems}
